If \(c_e(s)=c_{j^\star}(s)\) for every \(s\), then \(h_e(t)=y_{c_{j^\star}(s_t)}(t)\).  Since \(\rho_{j^\star}>0\), randomization and the common outcome response give, for every compatible \(v\),
\[
\mathbb E_v[Y\mid D=j^\star]
=\rho_{j^\star}^{-1}\sum_tv_t\rho_{j^\star}y_{c_{j^\star}(s_t)}(t)
=h_e^{\top}v.
\]
The left side is fixed by \(p\), so every \(v\) in the fiber has that target.  Sharpness gives \(L(p)=U(p)=\mathbb E[Y\mid D=j^\star]\).

If \(A_c=1\) for every \(c\), then \(R=1\) and \(Z_{\mathrm{obs}}=Z\) almost surely.  Nevertheless full revelation alone need not identify a distinct configuration.  Take one stratum, a trialed configuration \((1,q_\ell)\), an untrialed configuration \((1,q_h)\), and \(q_\ell<q_h\).  In two deterministic SCMs set \(Z=1\), all actions to one, and all neutral outcomes to zero.  Set
\[
(Y_{(1,q_\ell)},Y_{(1,q_h)})=(0,0)
\quad\text{in the first model},\qquad
(Y_{(1,q_\ell)},Y_{(1,q_h)})=(0,1)
\quad\text{in the second}.
\]
Both types satisfy correctness, neutrality, and both performance orders.  Their fully labeled trial laws coincide, but their untrialed means are zero and one.  This proves that revelation alone is insufficient while leaving intact the stated special-law, overlap, neutral-invariance, and further-restriction exceptions.